% !TEX root = ../Main.tex
\chapter{极限思想}

\section{无穷与趋近}

\textbf{无穷}记作 $\infty$，分两个方向：
\[
\infty:\quad
\bc
+\infty,\\
-\infty.
\ec
\]
$+\infty$ 是沿数轴向右越走越远、超过任何一个实数，$-\infty$ 则向左；二者合起来要把\textbf{全体实数}都包进去。

\textbf{趋近}记作 $\to$，表示一个变量越来越靠近某个目标值。先看三个最基本的例子：
\[
y_1=\frac{1}{x}:\quad x\to 0,\ y\to+\infty;\qquad x\to+\infty,\ y\to 0;
\]
\[
y_2=x:\quad x\to+\infty,\ y\to+\infty;\qquad
y_3=x^2:\quad x\to+\infty,\ y\to+\infty.
\]

$y_2$ 和 $y_3$ 当 $x\to+\infty$ 时都趋于 $+\infty$。那么是不是就有 $y_2=y_3$ 呢？\textbf{不是}。同样是 $+\infty$，跑得有快有慢——$y_3=x^2$ 比 $y_2=x$ 涨得快得多，记作
\[
y_3\gg y_2.
\]
所以"趋于无穷"只说明方向，还要比"谁涨得更快"。

\begin{center}
\begin{tikzpicture}[>=Stealth,scale=0.8]
  \draw[->] (-0.4,0) -- (3,0) node[right]{$x$};
  \draw[->] (0,-0.4) -- (0,4) node[above]{$y$};
  \draw[thick,blue!70!black,domain=0:1.9,smooth,samples=60] plot (\x,{\x*\x}) node[right]{\footnotesize$y_3=x^2$};
  \draw[thick,red!60!black,domain=0:2.8] plot (\x,{\x}) node[right]{\footnotesize$y_2=x$};
\end{tikzpicture}
\end{center}

\section{有理式在无穷处的极限}

下面四道题都是求 $x\to+\infty$ 时分式 $t$ 的极限。同一道题往往有好几种做法，目的是让你从不同角度都能看出"谁压倒谁"。

\ex{求 $\displaystyle t=\frac{x+1}{x+2}$ 当 $x\to+\infty$ 的极限。}

\sol{
    \textbf{分离常数（化为反比例）.} 把分子凑出分母：
    \[
    t=\frac{x+2-1}{x+2}=1-\frac{1}{x+2}.
    \]
    当 $x\to+\infty$ 时 $\dfrac{1}{x+2}\to 0$，故 $t\to 1$。
}

\ex{求 $\displaystyle t=\frac{x^2+x+3}{x+2}$ 当 $x\to+\infty$ 的极限。}

\sol{
    \textbf{法一（增长速率趋近）.} $x\to+\infty$ 时，分子由最高次项主导，分母也是：
    \[
    x^2+x+3=x^2\pr{1+\frac1x+\frac{3}{x^2}}\approx x^2,\qquad x+2\approx x,
    \]
    故
    \[
    t\approx\frac{x^2}{x}=x\ \to\ +\infty.
    \]

    \textbf{法二（常规换元）.} 令 $m=x+2$，即 $x=m-2$（$m\ne 0$），$x\to+\infty\Rightarrow m\to+\infty$：
    \[
    t=\frac{(m-2)^2+(m-2)+3}{m}=\frac{m^2-4m+4+m-2+3}{m}=\frac{m^2-3m+5}{m}=m-3+\frac{5}{m}.
    \]
    当 $m\to+\infty$，$m\to+\infty$、$\dfrac{5}{m}\to 0$，故 $t\to+\infty$。
}

\ex{求 $\displaystyle t=\frac{x+1}{x^2+2}$ 当 $x\to+\infty$ 的极限。}

\sol{
    \textbf{法一（换元）.} 令 $m=x+1$，即 $x=m-1$，则
    \[
    t=\frac{m}{(m-1)^2+2}=\frac{m}{m^2-2m+3}=\frac{1}{m-2+\frac{3}{m}}\ \to\ 0
    \quad(m\to+\infty).
    \]

    \textbf{法二（模型迁移）.} 取倒数：
    \[
    \frac{1}{t}=\frac{x^2+2}{x+1}.
    \]
    这正是上一题"二次比一次"的形式，当 $x\to+\infty$ 时 $\dfrac{1}{t}\to+\infty$，故 $t\to 0$。

    \textbf{法三（增长速率趋近）.} 只看最高次：
    \[
    t\approx\frac{x}{x^2}=\frac{1}{x}\ \to\ 0.
    \]
}

\ex{求 $\displaystyle t=\frac{2x^2+3x+1}{x^2+2x+3}$ 当 $x\to+\infty$ 的极限。}

\sol{
    \textbf{法一（分离常数）.} 凑出分母：
    \[
    t=\frac{2(x^2+2x+3)-x-5}{x^2+2x+3}=2-\frac{x+5}{x^2+2x+3}.
    \]
    余下的 $\dfrac{x+5}{x^2+2x+3}$ 正是上一题那种"低次比高次"的形式（换元或看增长速率都行），当 $x\to+\infty$ 时趋于 $0$，故 $t\to 2$。

    \textbf{法二（增长速率趋近）.} 直接比最高次：
    \[
    t\approx\frac{2x^2}{x^2}=2.
    \]
}
